\documentclass[11pt,a4paper]{article}
\usepackage[margin=2.5cm]{geometry}
\usepackage{amsmath,amssymb,amsthm,mathtools,booktabs,xcolor}
\usepackage[colorlinks=true,linkcolor=blue!60!black,citecolor=green!40!black]{hyperref}
\newcommand{\La}{\textsf{\textbf{(a)}}}\newcommand{\Lb}{\textsf{\textbf{(b)}}}
\newcommand{\Lc}{\textsf{\textbf{(c)}}}\newcommand{\Ld}{\textsf{\textbf{(d)}}}
\newcommand{\script}[1]{\texttt{corrigan-zambon/#1}}
\newtheorem{proposition}{Proposition}\theoremstyle{remark}\newtheorem{remark}{Remark}
\title{Semiclassical soliton--defect scattering\\ for the good Boussinesq equation}
\author{Project \emph{quantum\_boussinesq\_solitons}, Phase 2c}
\date{3 November 2026}
\begin{document}\maketitle
\begin{abstract}
We apply the symplectic (Darboux-coordinate) method of the Phase~2 report to solitons crossing the integrable type-I defect of Corrigan and Zambon (CZ) \cite{CZ}. Everything is for the good equation in the project convention. We find:
\begin{enumerate}
\item The CZ transmission factor is a ratio of KP roots, $z=[(p_b-q)/(p_a-q)]^{\eta}$. Here $(p_a,p_b)$ are the soliton's two roots and $q=-s/(2\sqrt3)$ is a single root carried by the defect.
\item The semiclassical transmission phase has the same structure as the two-soliton phase:
\[
\delta_D=\tfrac29\big(\ell-\ell(q)\big)\log z^{\epsilon}-\tfrac23\big(\Pi M_D-\Pi_D M\big).
\]
The defect's spectral parameter is $\ell(q)=(3s-s^3)/2$, and $(M_D,\Pi_D)$ are effective defect charges.
\item The type-I defect has no continuous degrees of freedom and adds no cross terms to $\Omega$. Two-soliton--defect scattering is symplectic with exactly the Phase~2 cross terms. It factorises: the total phase is $\delta_{12}+\delta_{1D}+\delta_{2D}$, whether the solitons collide before or after crossing the defect.
\item Absorption of a soliton by the defect happens exactly when $q$ coincides with one of its roots. The defect then moves to another of its vacua, and we verify exact energy--momentum balance. Phases stay finite on absorption and resonance lines, while delays diverge logarithmically.
\item Defect fusion gives an exact leading-order bootstrap, $\delta_D^{(q')}-\delta_D^{(q)}=\pm[\delta(3,S)+M_S\Pi_3]$. The bad-equation candidate S-matrix, continued to the good equation, fails the resulting triangle closure at $O(\hbar^0)$, by $e^{\pm i\pi/2}$; the pure exponential $e^{i\delta/\hbar}$ passes it exactly (Section~\ref{sec:cont}). However, a defect vacuum is linearly stable iff $\eta\sigma<-1$, so exactly one vacuum per triple is stable. The triangle argument therefore rests on an unstable intermediate state and is only heuristic (Section~\ref{sec:stab}).
\end{enumerate}
All checks are at 40 digits and are in \script{defect\_transmission\_phase.py} and \script{defect\_two\_soliton\_symplectic.py}; both end with \texttt{ALL CHECKS PASSED}.
\end{abstract}

\section{Set-up and dictionary}

The project convention is $u_{tt}=u_{xx}-u_{xxxx}-3(u^2)_{xx}$ with $u=\phi_x=2(\log\tau)_{xx}$. The Lagrangian is $L=\tfrac12\phi_t^2-\tfrac12\phi_x^2+\phi_x^3-\tfrac12\phi_{xx}^2$, the canonical form is $\Omega=\int d\pi\wedge d\phi$ with $\pi=\phi_t$, and $\hbar$ refers to this normalisation. A soliton has $k=\sin\alpha$, velocity $v=\epsilon\cos\alpha$, $M=\int u=2k$, $\Pi=\int\pi=-Mv$, $P=\tfrac23\epsilon k^3\cos\alpha$, and $\ell=-\epsilon\cos3\alpha$.

CZ \cite{CZ} (full text read) use a potential field $u_{\rm CZ}$ and set $c=\lambda=\kappa=1$. Translating:
\[
u_{\rm CZ}=-3\phi,\qquad L_{\rm CZ}=9L,\qquad P_{\rm CZ}=9P,\qquad \gamma=\operatorname{sech}\theta=k,\qquad \nu=\tanh\theta=\epsilon\cos\alpha .
\]
\Lc, check~1 of \script{defect\_transmission\_phase.py}.

Their defect has parameter $s=\sigma_0$ and orientation $\eta=\pm1$, where $\alpha_{\rm CZ}=\eta/\sqrt3$. Its sewing conditions were rederived from the action and checked on the exact solutions in \texttt{corrigan-zambon/verify\_defect\_solutions.py}. CZ's one- and two-soliton solutions are purely transmitting, with the factor
\[
z=\frac{\sinh\theta+s\cosh\theta-\eta\sqrt3}{\sinh\theta+s\cosh\theta+\eta\sqrt3}
\]
multiplying each $E_j$ on the far side \La\ (CZ (5.5), (5.11)).

\section{Symplectic structure with a type-I defect}

The defect term $D=u_tF+v_tG+H$ adds $F\,\delta u+G\,\delta v$ at $x_0$ to the canonical one-form. Asymptotically $u_x=v_x=0$ at $x_0$, so $F\to f$ and $G\to g$. Hence
\[
\Omega_D=(f_v-g_u)\,du\wedge dv\big|_{x_0}=\sigma\,du\wedge dv .
\]
The defect is asymptotically in a vacuum, where $\tau'(u-v)=0$, so $w=u-v$ at $x_0$ is a discrete label and $du=dv$. Therefore \emph{$\Omega_D=0$ on every vacuum sector}. \Lc, check~6.

It follows that:
\begin{itemize}
\item The type-I defect has no continuous degree of freedom.
\item There is no $d\Pi\wedge dM_D$ or $d\Pi_D\wedge dM$ cross term: both $M_D\propto w$ and the defect's shift charge are locally constant.
\end{itemize}
We use the shift-invariant gauge $f=g=-\sigma^2/(4\mu)$, which has $f_v-g_u=\sigma$. (CZ's (3.47) gives $2\sigma$; see \texttt{corrigan-zambon/README.md}.) In this gauge the total $\Pi$ is conserved through the defect.

For one soliton the phase space is $\Omega=dP\wedge dX+d\Pi_{\rm tot}\wedge dc$, with $X$ measured from the fixed defect. Any map $X\mapsto X+\Delta x(P)$ is then symplectic. One-soliton--defect scattering is therefore trivially symplectic, and the content lies in the phase and in the multi-soliton case.

\section{The transmission factor in KP roots}

\begin{proposition}[\Lc, check~2]\label{prop:z}
Let the soliton roots be $p_a=\tfrac12(k+\epsilon\cos\alpha/\sqrt3)$ and $p_b=\tfrac12(\epsilon\cos\alpha/\sqrt3-k)$, so that $p_a-p_b=k$ and $p_a^2+p_ap_b+p_b^2=\tfrac14$. Let $q=-s/(2\sqrt3)$. Then
\[
z=\Big(\frac{p_b-q}{p_a-q}\Big)^{\eta},\qquad
\ell=\ell(p_a)=\ell(p_b),\quad \ell(p):=-3\sqrt3\,p\,(1-4p^2).
\]
The two-soliton factor is the KP cross-ratio,
\[
A_{12}=S_{12}=\frac{(p_{1a}-p_{2a})(p_{1b}-p_{2b})}{(p_{1a}-p_{2b})(p_{1b}-p_{2a})} .
\]
(\Lc, check~1 of \script{defect\_two\_soliton\_symplectic.py}.) Formally, $z^{-\eta}$ is the cross-ratio of the soliton with a pair $(q,\infty)$. The defect behaves like ``half a soliton''.
\end{proposition}

CZ remark that no factor of $S_{12}$ seems closely related to $z$ \La. In root variables the relation is transparent.

\paragraph{Regular region.} $z>0$ iff $q\notin[p_b,p_a]$. The roots of all solitons lie in $[-1/\sqrt3,1/\sqrt3]$, so for $|s|>2$ every soliton is regular and none can be absorbed. This agrees with CZ's statement for $|s|>2$ \La. For $z<0$ the transmitted solution is singular \La.

\paragraph{Delay.} From peak positions of the exact solution, $\Delta x=-k^{-1}\log z^{\epsilon}$, where $z^\epsilon$ is the factor in the direction of motion \Lc\ (check~3, 16 regular cases, agreement to $10^{-39}$).

\section{The semiclassical transmission phase}

With $\Delta x=-\partial\delta_D/\partial P$ we have $\partial_\alpha\delta_D=\tfrac23\epsilon\sin3\alpha\,\log z^\epsilon$.

\begin{proposition}[\Lc, check~4]\label{prop:dD}
For all $\epsilon,\eta=\pm1$ and all $s$, in the regular region,
\begin{equation}\label{eq:dD}
\boxed{\ \delta_D=\tfrac29\big(\ell-\ell_D\big)\log z^{\epsilon}-\tfrac23\big(\Pi M_D-\Pi_D M\big)\ }
\end{equation}
with
\[
\ell_D=\ell(q)=\tfrac12(3s-s^3),\qquad M_D=\epsilon\eta\,q,\qquad \Pi_D=\epsilon\eta\,\frac{s^2-2}{2\sqrt3} .
\]
This holds up to an additive $s$-dependent constant, which the classical data do not fix. Moreover $\delta_D(-\eta)=-\delta_D(\eta)$: the two orientations are semiclassically inverse to each other.
\end{proposition}

\begin{remark}
\begin{itemize}
\item The coefficients $\ell_D,M_D,\Pi_D$ were first found by solving a linear system numerically. They were then identified exactly. For $\epsilon=\eta=1$, an exploratory polynomial division in $\cos\alpha$ (not part of the scripts) shows that the remainder has no pole at $z=0,\infty$ only for $\ell_D=(3s-s^3)/2$. The identification was then verified at 40 digits.
\item Compare the two-soliton phase $\delta_{12}=\tfrac29(\ell_1-\ell_2)\log A_{12}-\tfrac23(\Pi_1M_2-\Pi_2M_1)$. The defect enters exactly like a second particle, with spectral parameter $\ell(q)$ and charges $(M_D,\Pi_D)$. \Lc\ for the formula; the ``particle'' reading is \Ld.
\item Since $\rho_0=-(s/2\mu)(1-s^2/3)$ and $\mu=\eta/\sqrt3$ in CZ, we have $\ell_D=-\sqrt3\,\eta\,\rho_0$. The defect's spectral parameter is its constant $\rho_0$. \Lc\ (one line of algebra)
\item $M_D$ and $\Pi_D$ depend on $\epsilon$. They are therefore effective coefficients of a mutual-statistics term, not intrinsic charges of the defect. (In Section~2 the defect has none.)
\end{itemize}
\end{remark}

\section{Two solitons and a defect: symplecticity and factorisation}

CZ's two-soliton solution (5.11) multiplies $E_j$ by $z_j$ and leaves $S_{12}$ unchanged \La. We tested this as follows (\script{defect\_two\_soliton\_symplectic.py}).

\begin{proposition}[\Lc]\label{prop:2s}
\begin{enumerate}
\item \emph{Shifts} (check~2). From peak positions of the exact solution at $t=\pm1000$, $\Delta x_j=\Delta x_j^{\rm ss}+\Delta x_j^{D}$, with agreement to $<10^{-27}$. This was done for overtaking and head-on collisions, both orientations, and $s$ of both signs. The shifts are \emph{identical} whether the soliton--soliton collision takes place left or right of the defect.
\item \emph{Symplecticity} (check~3). The identity
\[
\sum_jdP_j\wedge d\Delta x_j=(d\Pi_i\wedge dM_j)_{\rm in}-(d\Pi_i\wedge dM_j)_{\rm out}
\]
holds to $10^{-40}$, with the Phase~2 cross terms and no defect cross terms. The left side alone is $O(1)$ (e.g.\ $-2.441$), so the naive test fails as in Phase~2.
\item \emph{Factorisation} (check~4). In the Phase~2 Darboux frame ($\phi$ normalised at $-\infty$),
\[
\Delta x'_j=-\partial_{P_j}\big(\delta_{12}+\delta_{1D}+\delta_{2D}\big),
\]
to $10^{-40}$.
\end{enumerate}
\end{proposition}

\paragraph{Interpretation.} Semiclassically the soliton--soliton--defect amplitude factorises as $S_{12}T_1T_2$. Because scattering is diagonal and purely transmitting, the reflection--transmission compatibility relation $T_1T_2S_{12}=S_{12}T_2T_1$ is automatic. The non-trivial content is twofold:
\begin{itemize}
\item the defect shifts are independent of the other soliton (an exact classical fact, \La/\Lc);
\item the Darboux-frame corrections $\gamma_jM_i$ are unaffected by the defect, because $\Omega_D=0$.
\end{itemize}
This is the semiclassical analogue of the sine-Gordon defect story, but with one difference. The required cross terms between solitons are present, while there are none between soliton and defect. \Lc\ for the leading order; nothing is claimed beyond it.

\section{Absorption, decay, the cusp and the resonances}

\subsection{Absorption and decay into one soliton}

The static defect vacua satisfy $\rho(\sigma)=0$ (which implies $\tau'=0$), that is $\sigma\in\{s,\,\sigma_\pm\}$ with $\sigma_\pm^2+s\sigma_\pm+s^2=3$. (An earlier version also listed $\sigma=0$. That is wrong: it solves $\tau'=0$ but violates the $u_{xx}$ sewing condition, Section~\ref{sec:stab}.) Write $\sigma=-2\sqrt3\,q'$. Then the condition for $\sigma_\pm$ reads
\[
q'^2+qq'+q^2=\tfrac14 .
\]
So the excited vacua correspond to the \emph{other two roots of the KP triple containing $q$}. \Lc

$z\in\{0,\infty\}$ (absorption, CZ \La) occurs iff $q=p_a$ or $q=p_b$. In that case we verify:
\[
(E,P)_{\rm soliton}=\pm\big[(E_D,P_D)(\sigma_\pm)-(E_D,P_D)(s)\big]\quad\text{exactly},
\]
with $E_D=-\tau$ and $P_D=h$ on vacua (CZ (3.48)--(3.51)). \Lc, check~5, $10^{-40}$.

In root language, the defect with root $q$ absorbs the soliton $(q,p)$ and becomes the defect with root $p$, in the same triple. The time-reversed process is the decay of an excited defect into a soliton plus the ground defect.

On an absorption line $\ell(p)=\ell(q)$, so the coefficient of $\log z$ in~\eqref{eq:dD} vanishes linearly. Hence:
\begin{itemize}
\item $\delta_D$ stays finite; numerically it converges to $-0.58063712\ldots$ at $s=0.7$;
\item the delay diverges like $\log|\alpha-\alpha^*|$.
\end{itemize}
\Lc

The natural quantum picture is that $T(\alpha)$ has a pole near $\alpha^*$, corresponding to a defect bound state $D_{q}\,S_{(q,p)}\to D_{p}$, as for sine-Gordon defects. \Ld\ This was not derived: a static excited vacuum is not a periodic orbit, so no Bohr--Sommerfeld rule applies directly.

\subsection{Resonance lines}

$A_{12}\to0$ or $\infty$ when two roots of different solitons coincide. For example, for overtaking $A\to0$ at $\alpha_1=2\pi/3-\alpha_2$, the mirror image of soliton~2 across the cusp, which has the same $\ell$ and the same root $p_a$. (An earlier version of this note said ``the same $P$''; that was wrong.) On these lines $\ell_1=\ell_2$, so
\[
\delta_{12}\to\tfrac23M_1M_2(v_1-v_2),
\]
which is finite. \Lc, check~5, for overtaking $A\to0$ and head-on $A\to\infty$. The delays diverge logarithmically. Past these lines lie the merging/splitting (Y-shaped) solutions, which are outside the two-particle description.

\subsection{The cusp $\alpha=\pi/3$}

$dP/d\alpha=\tfrac23k\sin3\alpha$ vanishes, and so does $d\delta_D/d\alpha$, while $\Delta x$ stays finite. In the example, $z=0.318$ or $8.5$ at $s=\pm2.4$. \Lc, check~5. So $\delta_D$ is smooth in $\alpha$ but has a square-root branch point as a function of $P$. The two branches are $\alpha<\pi/3$ and $\alpha>\pi/3$, and they share each value $P<P_{\max}$. The stationary-phase link $\Delta x=-\partial\delta/\partial P$ still holds on each branch. Near $P_{\max}$, however, the wave-packet derivation needs a uniform (Airy-type) treatment, and $\alpha$ (or the root) rather than $P$ is the natural label. \Ld

We therefore state results in the regular region: $\alpha_j<\pi/3$ for the two-soliton checks, $z_j>0$, and $A_{12}>0$.

\subsection{Status of reflection and of decay into two solitons}

\begin{itemize}
\item For ordinary one- and two-soliton solutions, the CZ defect is \emph{purely transmitting} \La. Hence $R=0$ at the semiclassical order considered here \Lc.
\item CZ's reflection occurs only for the merging/splitting three-exponential solutions (their \S5.3) \La. Those lie on or beyond the resonance lines and involve soliton-number change. The phase-space description used here (fixed number of asymptotic solitons, Darboux coordinates) does not cover them. We have \emph{not} derived a reflection amplitude.
\item Likewise, CZ's ``decaying defect'' into two solitons (their Fig.~8) has not been analysed.
\item Our guess is that semiclassical reflection is a resonance phenomenon: reflection $\sim$ absorption of the incoming soliton into an excited vacuum, followed by emission into the other root pair of the triple. Quantum reflection amplitudes would then be suppressed except near the resonance and absorption lines. \Ld
\end{itemize}

\section{Test of the continued bad-equation S-matrix}\label{sec:cont}

\paragraph{Set-up.} We take the candidate of \texttt{smatrix\_ansatz.pdf} with good-equation data: $\ell=-\epsilon\cos3\alpha$, $A_{12}$ the good cross-ratio, and $\eta=\tfrac29(\ell_1-\ell_2)/\hbar$. With $R(x)=\Gamma(1+ix)/\Gamma(1-ix)$,
\[
S_{\epsilon\epsilon}=R(\eta)\Big(\frac{e^2A_{12}}{\eta^2}\Big)^{i\eta}e^{iT_{12}/\hbar},\qquad
S_{+-}=R(\eta)^{-1}\Big(\frac{A_{12}\eta^2}{e^2}\Big)^{i\eta}e^{iT_{12}/\hbar}.
\]
The good-equation version of the semiclassical phase is
\[
\delta=\tfrac29(\ell_1-\ell_2)\log A_{12}+T_{12}.
\]
The candidate is \Ld; everything below is \Lc, from \texttt{phase2/scripts/defect\_smatrix\_good\_continuation.py}.

\paragraph{Regular region.} For $\alpha_j<\pi/3$ and $A_{12}>0$ the continued candidate passes the same tests as the bad one (check~1):
\begin{itemize}
\item $|S|=1$ and $S(1,2)S(2,1)=1$;
\item $\arg S-\delta/\hbar\to\tfrac\pi2\tau\,\mathrm{sgn}\,\eta$, with remainder $O(\hbar)$; for example $5.5\times10^{-3},\,5.5\times10^{-4},\,5.5\times10^{-5}$ at $\hbar=10^{-2},10^{-3},10^{-4}$.
\end{itemize}

\paragraph{Leading-order defect fusion.} A defect with root $q$ absorbs the soliton $S=\{q,q'\}$ and becomes the defect with root $q'$ (Section~6). For every test soliton~3 in the regular region we find, exactly,
\begin{equation}\label{eq:fusion}
\delta_D^{(q')}(3)-\delta_D^{(q)}(3)=\rho\,\big[\delta(3,S)+M_S\Pi_3\big],\qquad\rho=\pm1 .
\end{equation}
\begin{itemize}
\item There is no free constant: the constants implicit in~\eqref{eq:dD} are consistent with zero.
\item The term $M_S\Pi_3$ is the cross term of soliton~3 with the $\phi$-jump $M_S$ stored in the excited defect.
\item Verified to $10^{-40}$ for 12 triangles (check~4).
\end{itemize}
This is the semiclassical soliton--soliton--defect bootstrap. Note that all three solitons of one KP triple have the same $\ell_S$.

\paragraph{Exact triangle closure.} Assume an exact bootstrap $T^{(q')}(3)=T^{(q)}(3)\,[S(3,S)e^{iM_S\Pi_3/\hbar}]^{\rho}$ \Ld. Going around the triangle $q\to q'\to q''\to q$, the defect amplitudes cancel. This leaves a condition on $S$ alone:
\[
\prod_{i=1}^3\big[S(3,S_i)\,e^{iM_i\Pi_3/\hbar}\big]^{\rho_i}=1 .
\]
For the continued candidate, every factor has the same $\eta=\tfrac29(\ell_3-\ell_S)/\hbar$. The powers of $A$ cancel by~\eqref{eq:fusion}, so the product equals
\[
\big[R(\eta)(e/\eta)^{2i\eta}\big]^{\sum_i\rho_i\tau_i},
\]
where $\tau_i=+1$ ($-1$) for same (opposite) direction. In all 12 triangles $\sum_i\rho_i\tau_i=\pm1$. The closure therefore fails by
\[
e^{\pm i\frac\pi2\mathrm{sgn}\,\eta}\big(1+O(\hbar)\big),
\]
for example closure phases $-1.4989,\,-1.5637,\,-1.5701$ at $\hbar=0.1,0.01,0.001$. \Lc

The failure:
\begin{itemize}
\item is invisible at $O(\hbar^{-1})$;
\item appears at $O(\hbar^0)$ as a step of $\pi$ when $\ell_3$ crosses $\ell_S$, so it is not a constant convention;
\item continues at all higher orders, through the Stirling series of $R$.
\end{itemize}
The pure exponential $S=e^{i\delta/\hbar}$, $T^{(q)}=e^{i\delta_D/\hbar}$ satisfies the closure \emph{exactly} ($10^{-38}$). \Lc

\paragraph{Resonance lines.} On the overtaking line $\alpha_1=2\pi/3-\alpha_2$ we have $\ell_1=\ell_2$, so $\eta\to0$, and $S\to e^{iT_{12}/\hbar}$ continuously; for example $|\log S-iT/\hbar|=1.7\times10^{-7}$ at distance $10^{-9}$. Beyond the line $A<0$ and $|S|\neq1$; for example $|S|=1.037$ at distance $10^{-3}$, $\hbar=0.05$. The candidate is not unitary in the singular/resonant region. \Lc (check~2)

\paragraph{Cusp.} $\ell$ is stationary at $\alpha=\pi/3$. Hence $\eta$ varies quadratically there, and the line $\eta=0$ through the cusp is the resonance line itself. The candidate is smooth in $\alpha$ through the cusp. The issue there is the labelling by $P$ (Section~6.3), not $S$. \Lc

\paragraph{Breathers.} At $\alpha_{1,2}=a\mp ib$ the Gamma poles $i\eta=-N$ exist. Examples:
\begin{itemize}
\item $a=0.4$, $\hbar=0.1$: $b=0.0797$ for $N=1$, with real total $P$ and $\operatorname{Im}(P_1-P_2)>0$;
\item at all four points tested, $A<0$ there.
\end{itemize}
Every Hirota breather of the good equation is singular (Phase~1, Prop.~\emph{nobreather}). So the continued candidate predicts bound states with no regular classical counterpart. \Lc\ for the pole positions; whether such quantum bound states exist is open.

\paragraph{Implications.} \Ld
\begin{enumerate}
\item In the bad equation, the Gamma factor was \emph{required} by the breather poles at the Bogdan--Charkina points. In the good equation that reason is absent: there are no regular breathers.
\item Defect fusion (if the exact bootstrap holds) actively disfavours the Gamma factor. A good-equation S-matrix that is compatible with integrable defects must have a Gamma part that is a \emph{coboundary} over KP triples. Since $\eta$ is the same within a triple, this holds iff the Gamma exponents satisfy $\sum_i\rho_i\tau_i=0$, which the continued form violates.
\item The simplest consistent choice is the pure exponential $e^{i\delta/\hbar}$, which was excluded for the bad equation \emph{only} by its lack of poles. Alternatively, a CDD-type factor $\kappa(\eta)$ with $\kappa^{\sum\rho\tau}$ cancelling $R(\eta)(e/\eta)^{2i\eta}$ would also work. These two options cannot be distinguished at the order tested.
\item So the good-equation S-matrix is \emph{not} the naive continuation $\beta\to i\alpha$ of the bad candidate. They can agree at $O(\hbar^{-1})$, and they must differ at $O(\hbar^0)$ (Maslov-type phase) and in the pole structure.
\item Caveats:
  \begin{itemize}
  \item The closure argument assumes that the defect bound states (excited vacua) are genuine asymptotic states, so that the bootstrap applies.
  \item Absorption happens at real rapidity, so these states are embedded in the continuum (the defect breaks $P$-conservation). Their quantum stability is not established.
  \item If they are resonances rather than stable states, the triangle condition need not hold exactly.
  \end{itemize}
\end{enumerate}

\section{Linear stability of the defect vacua}\label{sec:stab}

The exact-bootstrap argument of Section~\ref{sec:cont} assumes that excited defect vacua are stable asymptotic states. We tested this classically, by linearising the sewing conditions about each static vacuum. The bulk equation becomes $a_{tt}=a_{xx}-a_{xxxx}$. Everything here is \Lc, from \script{defect\_vacuum\_stability.py}.

\begin{proposition}[\Lc]\label{prop:stab}
\begin{enumerate}
\item \emph{Vacua.} The static vacua are exactly $\sigma\in\{s,\sigma_\pm\}$ (check~1).
\item \emph{Phonons.} On every real vacuum, stable or not, phonons are purely transmitted: $R=0$ and $|T|=1$. Moreover $T$ equals CZ's linear-defect formula~(4.6) with $\sigma_0\to\sigma_{\rm vac}$ and $\alpha+\beta=\sqrt3\eta$:
\[
T_{\rm ph}=-\frac{\sqrt3\eta k^2+i\sigma k+i\omega}{\sqrt3\eta k^2-i\sigma k-i\omega},
\]
to $10^{-30}$ (check~2).
\item \emph{Localised modes.} The spectrum depends on $(\eta,\sigma_{\rm vac})$ only, not on $s$. A vacuum has a localised growing mode ($\operatorname{Im}\omega>0$) iff
\[
\eta\,\sigma_{\rm vac}>-1 .
\]
\begin{itemize}
\item The growth rate tends to $0$ at $\eta\sigma=-1$; for example $\gamma=0.029$ at $\eta\sigma=-0.95$ and $\gamma=0.4999$ at $\eta\sigma=0.5$.
\item A second mode appears for $\eta\sigma>1$, and a complex pair for $\eta\sigma>2$.
\item Every growing mode on the imaginary axis has zero conserved energy, to $10^{-20}$ relative. This is a necessary condition for growth and an independent check of the linearisation. The defect's quadratic energy $-H_2$ is indefinite, which is what allows the instability.
\end{itemize}
(Check~3. The root search is a fine scan of the imaginary axis plus Newton iteration from a complex grid. It is numerical, not a proof of completeness. The branch point $\omega=i/2$, where the two decay rates coincide, gives a spurious zero and is excluded.)
\item \emph{One stable vacuum per triple.} The vacua of one defect are $\sigma=-2\sqrt3\,q'$ for the three roots of one KP triple. Exactly one root has $\eta q'>1/(2\sqrt3)$, so \emph{exactly one vacuum per triple is linearly stable} (check~4).
\end{enumerate}
\end{proposition}

\paragraph{Consequences.}
\begin{itemize}
\item A CZ defect with ground vacuum $\sigma=s$ is linearly stable iff $\eta s<-1$.
\item CZ's illustrative example ($\theta=1$, $s=0.5$, and $z=0.05836$, which requires $\eta=+1$) has a linearly unstable ground vacuum, with growth rate $\approx0.4999$. Its exact soliton solutions are unaffected, but generic small perturbations of the defect grow. \Lc. The nonlinear fate of the instability was not studied.
\item For a stable ground defect ($\eta s<-1$), \emph{both} excited vacua are unstable. Absorption is possible only for $-2<\eta s<-1$, and the absorbed soliton leaves the defect in an unstable state. That state is presumably a resonance that re-emits rather than a bound state. \Ld
\item The hypothesis of the exact triangle closure in Section~\ref{sec:cont} therefore fails at the linear level. The $O(\hbar^0)$ objection to the continued Gamma candidate is downgraded from a conditional \Lc\ statement to a heuristic \Ld. The leading-order fusion identity~\eqref{eq:fusion} is unaffected: it is a classical identity.
\item The semiclassical transmission phase~\eqref{eq:dD}, symplecticity and factorisation are classical scattering statements and remain valid. A quantum defect requires a stable vacuum, so they are physically relevant for $\eta s<-1$.
\item The reflectionless phonon transmission $T_{\rm ph}$ is exactly what a one-loop (DHN-type) treatment of soliton--defect scattering needs.
\end{itemize}

\paragraph{Next step.} The $O(\hbar^0)$ question (the Maslov-type $\pm\pi/2$ of the Gamma candidate versus the pure exponential) must be decided directly by the one-loop two-soliton phase of the good equation. The defect bootstrap no longer decides it.

\section{Summary of labels and checks}

\begin{center}\begin{tabular}{lll}\toprule
Claim & Label & Check\\\midrule
Root form of $z$ and of $S_{12}$ & \Lc & phase script 2; two-soliton script 1\\
Delay from the exact solution & \Lc & phase script 3\\
$\delta_D$, eq.~\eqref{eq:dD} & \Lc & phase script 4\\
Vacua $\leftrightarrow$ triple; absorption energy balance & \Lc & phase script 5\\
$\Omega_D=0$ on vacua, no defect cross terms & \Lc & phase script 6\\
2-soliton shifts, order independence & \Lc & two-soliton script 2\\
Symplecticity with Phase 2 cross terms & \Lc & two-soliton script 3\\
Factorisation $\delta_{12}+\delta_{1D}+\delta_{2D}$ & \Lc & two-soliton script 4\\
Finite phases on resonance/absorption lines; cusp & \Lc & two-soliton script 5\\
Bound-state poles, reflection as resonance, cusp uniformisation & \Ld & ---\\
Fusion~\eqref{eq:fusion}; continued candidate fails triangle closure by $e^{\pm i\pi/2}$ & \Lc & \texttt{defect\_smatrix\_good\_continuation.py}\\
Good S-matrix must differ from continued bad one at $O(\hbar^0)$ & \Ld & ---\\
Vacua; phonon $T$; stable iff $\eta\sigma<-1$; one stable vacuum per triple & \Lc & \texttt{defect\_vacuum\_stability.py}\\\bottomrule
\end{tabular}\end{center}

\paragraph{Failures recorded.}
\begin{itemize}
\item A first version of the delay check started the root-finder too far from the peak and returned wrong roots. It was fixed by starting near the predicted peak and requiring the peak height $k^2/2$.
\item A first resonance test converged to a spurious root. It was replaced by the explicit line $p_{1a}=p_{2a}$ (or $p_{1b}=p_{2a}$).
\item For $s=1.9$ only one absorbable soliton exists; the check originally demanded two.
\end{itemize}

\begin{thebibliography}{9}
\bibitem{CZ} E.~Corrigan, C.~Zambon, \emph{Adding integrable defects to the Boussinesq equation}, J.~Phys.~A 56 (2023), arXiv:2306.03772. [full text read]
\end{thebibliography}
\end{document}
